
% =============================================
% CONFIGURACAO DO DOCUMENTO
% =============================================

\usepackage[utf8]{inputenc}
\usepackage[T1]{fontenc}
\usepackage[portuguese]{babel}
\usepackage{lmodern}
\usepackage{microtype}

\usepackage{geometry}
\geometry{
    a4paper,
    left=2.3cm,
    right=2.3cm,
    top=2.5cm,
    bottom=2.3cm
}

\usepackage[table]{xcolor}

% Cores
\definecolor{azulPDMI}{HTML}{1D5787}
\definecolor{azulClaro}{HTML}{EAF3FA}
\definecolor{cinzaLinha}{HTML}{D0D8DF}

\colorlet{azul}{azulPDMI}
\colorlet{azulclaro}{azulClaro}
\definecolor{verde}{HTML}{248563}

% Pacotes matematicos e graficos
\usepackage{amsmath,amssymb}
\usepackage{array,tabularx,longtable}
\usepackage{tikz}
\usetikzlibrary{arrows.meta}
\usepackage[most]{tcolorbox}

% Organizacao e formatacao
\usepackage{fancyhdr}
\usepackage{enumitem}
\usepackage{titlesec}
\usepackage{lastpage}
\usepackage{needspace}

% Espacamento geral
\setlength{\parindent}{0pt}
\setlength{\parskip}{4pt}
\setlength{\headheight}{26pt}
\renewcommand{\arraystretch}{1.55}

% =============================================
% CABECALHOS E RODAPES
% =============================================

\pagestyle{fancy}
\fancyhf{}

\newcommand{\currentparttitle}{}

\newcommand{\partwithheader}[1]{%
  \part{#1}%
  \gdef\currentparttitle{Parte \thepart\ -- #1}%
  \markboth{\currentparttitle}{}%
}

\renewcommand{\partmark}[1]{%
  \gdef\currentparttitle{Parte \thepart\ -- #1}%
  \markboth{\currentparttitle}{}%
}

\renewcommand{\chaptermark}[1]{%
  \markboth{\currentparttitle}
  {\chaptername\ \thechapter\ -- #1}%
}

\fancyhead[L]{%
  \small\textcolor{azulPDMI}{%
    \textbf{\nouppercase{\leftmark}}%
  }%
}

\fancyhead[R]{%
  \small\textcolor{azulPDMI}{%
    \textbf{\nouppercase{\rightmark}}%
  }%
}

\fancyfoot[R]{%
  \footnotesize Página \thepage\ de \pageref{LastPage}
}

\renewcommand{\headrulewidth}{0.4pt}
\renewcommand{\footrulewidth}{0pt}

\renewcommand{\headrule}{%
  \hbox to\headwidth{%
    \color{azulPDMI!35}%
    \leaders\hrule height \headrulewidth\hfill
  }%
}

\fancypagestyle{plain}{%
  \fancyhf{}
  \fancyhead[L]{%
    \small\textcolor{azulPDMI}{%
      \textbf{\nouppercase{\leftmark}}%
    }%
  }
  \fancyhead[R]{%
    \small\textcolor{azulPDMI}{%
      \textbf{\nouppercase{\rightmark}}%
    }%
  }
  \fancyfoot[R]{%
    \footnotesize Página \thepage\ de \pageref{LastPage}
  }
  \renewcommand{\headrulewidth}{0.4pt}
  \renewcommand{\footrulewidth}{0pt}
}

% =============================================
% FORMATACAO DOS CAPITULOS
% =============================================

\titleformat{\chapter}[display]
  {\normalfont\bfseries}
  {\large\chaptername\ \thechapter}
  {6pt}
  {\LARGE}

% Reduz o espaco acima e abaixo dos titulos
\titlespacing*{\chapter}
  {0pt}
  {3pt}
  {14pt}

% =============================================
% FORMATACAO DAS SECCOES
% =============================================

\titleformat{\section}
  {\large\bfseries\color{azulPDMI}}
  {}
  {0pt}
  {}

\titlespacing*{\section}
  {0pt}
  {10pt}
  {5pt}

% =============================================
% COMANDOS PERSONALIZADOS
% =============================================

\newcommand{\linhas}[1]{%
  \par\vspace{2pt}%
  \begingroup\color{cinzaLinha}%
  \foreach \n in {1,...,#1}{%
    \par\noindent\rule{\linewidth}{0.45pt}%
    \vspace{0.36cm}%
  }%
  \par\endgroup\vspace{2pt}%
}

\newcommand{\pergunta}[1]{%
  \textbf{#1}\enspace
}

\newcommand{\parte}[2]{%
  \section*{#1%
    \if\relax\detokenize{#2}\relax
    \else\ --- #2\fi
  }%
}

% =============================================
% CAIXAS
% =============================================

\newtcolorbox{desafio}{
  colback=azulClaro,
  colframe=azulPDMI!45,
  boxrule=0.5pt,
  arc=2mm,
  left=8pt,
  right=8pt,
  top=6pt,
  bottom=6pt
}

\newtcolorbox{nota}{
  colback=azulClaro,
  colframe=azulPDMI!35,
  boxrule=0.4pt,
  arc=1mm,
  left=8pt,
  right=8pt,
  top=6pt,
  bottom=6pt
}

\newtcolorbox{conclusao}{
  colback=azulClaro,
  colframe=azulPDMI!45,
  boxrule=0.5pt,
  arc=2mm,
  left=8pt,
  right=8pt,
  top=6pt,
  bottom=6pt
}

% =============================================
% POLIGONOS DA TABELA DA AULA 1
% =============================================

\newcommand{\poligonoTabela}[1]{%
  \begin{tikzpicture}[
    x=0.83cm,y=0.72cm,
    baseline=(current bounding box.center),
    line join=round
  ]
    \ifcase#1\relax
    \or\or\or
      \draw[draw=azulPDMI,fill=azulClaro,thick]
        (0.15,0.15)--(3.30,0.12)--(1.65,2.10)--cycle;
    \or
      \draw[draw=azulPDMI,fill=azulClaro,thick]
        (0.15,0.20)--(3.30,0.15)--(3.20,2.15)--(0.40,2.30)--cycle;
    \or
      \draw[draw=azulPDMI,fill=azulClaro,thick]
        (0.15,0.55)--(1.20,0.08)--(3.10,0.40)--
        (3.15,1.85)--(1.45,2.35)--cycle;
    \or
      \draw[draw=azulPDMI,fill=azulClaro,thick]
        (0.30,0.55)--(1.10,0.07)--(2.55,0.15)--
        (3.35,1.10)--(2.75,2.25)--(0.65,2.15)--cycle;
    \or
      \draw[draw=azulPDMI,fill=azulClaro,thick]
        (0.15,0.85)--(0.55,0.12)--(1.85,0.08)--
        (3.20,0.75)--(3.05,1.75)--(2.05,2.35)--
        (0.55,2.20)--cycle;
    \fi
  \end{tikzpicture}%
}

% Designacao dos capitulos
\addto\captionsportuguese{%
  \renewcommand{\chaptername}{Secção}%
}
